\section{Introducción}

 El desarrollo de sistemas modernos exige un dominio sólido sobre cómo se almacenan, organizan y recuperan los datos. Mediante el uso del lenguaje SQL y el gestor de base de datos PostgreSQL, es posible diseñar estructuras de información eficientes que respondan directamente a las necesidades del negocio o del proyecto en cuestión.

Este trabajo tiene como propósito plasmar la comprensión y desarrollo de una base de datos mediante SQL. Para ello, se presenta el análisis de los requerimientos asignados, la creación de esquemas relacionales y la aplicación de sentencias DDL (Data Definition Language) y DML (Data Manipulation Language) necesarias para construir una solución funcional y estructurada.
